\scanpage{67}\mspage{B.}
on peut donc supposer que
$Z$ est \uline{projectif}, \deleted{Si $\dim$} et $=$ \uline{normal}
si on veut. Si \inserted{$\dim Z\leq 0$, trivial. Si} $\dim Z=1$, on aura $Y=\varnothing$,
i.e. $X=Z$, et le th. est connu. Si $\dim Z\geq 2$
on utilise ceci :

\uline{Lemme} de \uline{Bertini} : (Zariski) Soit $k$ corps,
$X\xrightarrow{f}P^r_k$ un \uline{morphisme}, avec $X$ \deleted{\unsure{intègre}}\inserted{irréductible},
$\dim f(X)=n$, soit $1\leq\unsure{s|5}\leq\unsure{n|r}-1$, \deleted{et} soit
$u$ \deleted{le} point générique de $\underbrace{P^r_k\times\cdots\times P^r_k}_{s}$, et
$K$ une clôture algébrique de $k(u)$, enfin
$L_u$ la variété linéaire dans $P^r_K$ définie
par $u$, et $X_u=X\times_{P^r_k}L_u=X_K\times_{P^r_K} L_u$.
\margin{peut-on remplacer ``irréductible'' par ``connexe'' ? \uline{non}}
% slanted
Alors $X_u$ est irréductible.

% Diagram in the left margin, spanning the lemma and corollary.
\diagram{\begin{tikzcd}[ampersand replacement=\&,row sep=small,column sep=small]
X'_u\arrow[d]\&{}\&X'\arrow[d]\\
X_u\arrow[r]\arrow[d]\&X_K\arrow[r]\arrow[d]\&X\arrow[d]\\
L_u\arrow[r]\&P^r_K\arrow[r]\&P^r_k
\end{tikzcd}}

\uline{Corollaire} \deleted{\unsure{$\pi_1(X)$|Il est}} Si $X$ est \uline{normal},
alors $\pi_1(X_u)\to\pi_1(X)$ est \uline{surjectif}.

En effet, si $X'$ est un revêtement \deleted{\unsure{normal|connexe}}
étale connexe de $X$, étant normal il est
\uline{irréductible}, donc $X'_u=X'\times_X X_u$ est
irréductible, donc connexe, OK.
\margin{Si on admet le th. de connexion, on \deleted{peut} \deleted{\unsure{peut}} \unsure{aussi} utiliser une section \deleted{\unsure{plus|fibrée}} ``spéciale'' ..... p. ex une courbe ...}
% slanted
% archive sheet number: 66
